\documentclass[11pt]{article}

\usepackage{math_article}
\usepackage{series_ra}

\renewcommand{\ArticleNumeral}{I}
\renewcommand{\ArticleTitle}{Measure Theory and\\Lebesgue Integration}
\renewcommand{\ArticleAuthor}{Anqiao Ouyang}
\renewcommand{\ArticleDate}{August 2026}

\begin{document}

\maketitle

\section*{Introduction}

This article introduces the foundations of measure theory and the Lebesgue integral, which together form the backbone of modern real analysis. We begin with $\sigma$-algebras and measures, construct the Lebesgue measure on $\mathbb{R}^n$, and develop the integral with respect to a general measure.

\section{Sigma-Algebras and Measurable Spaces}

\begin{definition}[$\sigma$-Algebra]
Let $X$ be a non-empty set. A collection $\mathcal{F} \subseteq \mathcal{P}(X)$ is a \emph{$\sigma$-algebra} on $X$ if:
\begin{enumerate}
  \item $X \in \mathcal{F}$;
  \item if $A \in \mathcal{F}$, then $A^c \in \mathcal{F}$;
  \item if $(A_n)_{n=1}^\infty$ is a sequence in $\mathcal{F}$, then $\bigcup_{n=1}^\infty A_n \in \mathcal{F}$.
\end{enumerate}
The pair $(X, \mathcal{F})$ is called a \emph{measurable space}.
\end{definition}

\begin{example}
The Borel $\sigma$-algebra $\mathcal{B}(\mathbb{R})$ is the $\sigma$-algebra generated by the open sets of $\mathbb{R}$. It contains all open sets, closed sets, countable intersections of open sets ($G_\delta$ sets), and countable unions of closed sets ($F_\sigma$ sets).
\end{example}

\section{Measures}

\begin{definition}[Measure]
Let $(X, \mathcal{F})$ be a measurable space. A function $\mu \colon \mathcal{F} \to [0, \infty]$ is a \emph{measure} if:
\begin{enumerate}
  \item $\mu(\varnothing) = 0$;
  \item (countable additivity) for any sequence $(A_n)$ of pairwise disjoint sets in $\mathcal{F}$,
  \[
    \mu\!\left(\bigcup_{n=1}^\infty A_n\right) = \sum_{n=1}^\infty \mu(A_n).
  \]
\end{enumerate}
\end{definition}

\begin{theorem}[Carath\'eodory Extension]
Let $\mu_0$ be a pre-measure on an algebra $\mathcal{A} \subseteq \mathcal{P}(X)$. Then there exists a measure $\mu$ on $\sigma(\mathcal{A})$ that extends $\mu_0$. If $\mu_0$ is $\sigma$-finite, the extension is unique.
\end{theorem}

\begin{proof}
Define the outer measure $\mu^*(E) = \inf \sum_{n} \mu_0(A_n)$, where the infimum is taken over all countable covers $\{A_n\} \subseteq \mathcal{A}$ of $E$. One then shows that the collection of $\mu^*$-measurable sets forms a $\sigma$-algebra containing $\mathcal{A}$, and $\mu^*$ restricted to this $\sigma$-algebra is a measure extending $\mu_0$.
\end{proof}

\section{The Lebesgue Integral}

\begin{definition}[Lebesgue Integral of a Simple Function]
Let $(X, \mathcal{F}, \mu)$ be a measure space and $\varphi = \sum_{i=1}^n a_i \mathbf{1}_{A_i}$ a non-negative simple function. The Lebesgue integral of $\varphi$ is
\[
  \int_X \varphi \, d\mu = \sum_{i=1}^n a_i \, \mu(A_i).
\]
\end{definition}

\begin{theorem}[Monotone Convergence Theorem]
Let $(f_n)$ be a sequence of non-negative measurable functions with $f_n \uparrow f$ pointwise. Then
\[
  \lim_{n \to \infty} \int_X f_n \, d\mu = \int_X f \, d\mu.
\]
\end{theorem}

\begin{intuition}
The Monotone Convergence Theorem allows us to exchange limits and integrals for increasing sequences. This is the engine that powers most constructions of the Lebesgue integral: approximate from below by simple functions and pass to the limit.
\end{intuition}

\begin{theorem}[Dominated Convergence Theorem]
Let $(f_n)$ be a sequence of measurable functions converging pointwise to $f$. If there exists an integrable function $g$ such that $|f_n| \leq g$ for all $n$, then
\[
  \lim_{n \to \infty} \int_X f_n \, d\mu = \int_X f \, d\mu.
\]
\end{theorem}

\begin{remark}
The domination condition $|f_n| \leq g$ with $g$ integrable cannot be dropped. The sequence $f_n = n \cdot \mathbf{1}_{(0, 1/n)}$ converges pointwise to $0$ on $(0,1)$, yet $\int f_n = 1$ for all $n$.
\end{remark}

\end{document}
